\section{Algorithm: Multi-scale MCMC Option B}
\label{sec:algorithm}

\subsection{Setup and Notation}

Let $G = (V, E)$ be the census-tract adjacency graph for a state, with
$|V| = n$ tracts.
Each tract $v$ has population $\mathrm{pop}(v)$ and a GEOID string
$\mathrm{geoid}(v)$ of 11 digits (2 state + 3 county + 6 tract).
The county FIPS code of tract $v$ is
\[
  \mathrm{county}(v) = \mathrm{geoid}(v)[0{:}5],
\]
the first five characters of the GEOID.
Let $C = \{\mathrm{county}(v) : v \in V\}$ be the set of distinct county
FIPS codes, $|C| = n_C$.

The ideal fine population is $P_{\mathrm{ideal}} = \sum_v \mathrm{pop}(v) / k$;
the fine tolerance is $\epsilon_f$ (default 0.005).
The coarse tolerance is $\epsilon_c = 3 \epsilon_f$ (default 0.015), to
reflect the lower precision of county-level population accounting.
The factor of 3 is chosen as follows: a county-level coarse move moves all $m_c \approx n/100$ tracts in a county simultaneously (for a 100-county state). The fine-level rebalance step then corrects any violations. The factor 3 ensures that for states where the largest single county holds up to $3\varepsilon \cdot \text{pop}(\text{state})$ of total population (e.g., Los Angeles County, CA), the coarse chain still proposes valid county-level moves rather than immediately rejecting. For states with more uniform county populations (e.g., Iowa), a factor of 2 would suffice; the factor 3 is a conservative default. \textbf{Practitioners in states with extreme county population concentration should increase this factor or switch to Option A (BG-level fine moves).}

\subsection{County Adjacency Derivation}

\begin{proposition}[County adjacency from tract graph]
\label{prop:county-adj-exact}
Define the county adjacency graph $G_C = (C, E_C)$ by
\[
  (c_1, c_2) \in E_C
  \iff c_1 \neq c_2
  \text{ and }
  \exists\, (v_1, v_2) \in E \text{ with }
  \mathrm{county}(v_1) = c_1,\;
  \mathrm{county}(v_2) = c_2.
\]
This derivation is exact: two counties are geographically adjacent if and
only if they share a border, which implies the existence of at least one
pair of adjacent tracts straddling the county boundary.
The derivation requires one pass over the tract edge list and runs in $O(|E|)$
time.
\end{proposition}

\begin{proof}
The implication ``geographic adjacency $\Rightarrow$ cross-county tract edge''
holds because county boundaries in the census TIGER/Line products are
composed of sequences of tract boundaries.
Every county boundary segment that forms a shared border between two counties
is simultaneously a boundary segment of at least one pair of tracts, one in
each county.
The reverse implication ``cross-county tract edge $\Rightarrow$ geographic
adjacency'' holds by definition: two tracts can share an edge in the
adjacency graph only if they share a boundary segment, and if the tracts
belong to different counties then those counties share that boundary segment.
\end{proof}

\begin{remark}[Water-boundary caveat]
Proposition~3.1 assumes that two geographically adjacent counties share at least one cross-county tract edge in the TIGER/Line adjacency graph. This holds for land-boundary adjacency. However, some counties are geographically adjacent only across water bodies (e.g., two counties on opposite banks of a river whose tract boundaries do not cross). In TIGER/Line adjacency graphs built without water-crossing bridges, such county pairs may be absent from the derived county adjacency graph---producing a county graph that is slightly sparser than geographic adjacency.

For redistricting purposes, this is typically acceptable: federal congressional districts are not required to cross water bodies, and the missing county adjacency represents a pair that would be difficult to connect in practice. Users in states with significant water-separated county pairs (e.g., Texas Gulf Coast, Louisiana bayou regions) should verify the derived county adjacency against geographic county maps.
\end{remark}

The derivation scans each edge $(v_1, v_2) \in E$; if
$\mathrm{county}(v_1) \neq \mathrm{county}(v_2)$, the pair is added to
$E_C$.
For Texas, $|E| \approx 15{,}600$ tract edges yield $|E_C| \approx 220$
county edges from $|C| = 254$ counties.

\subsection{County Population and Initial Coarse Plan}

The county population is
$\mathrm{pop}(c) = \sum_{v \in f^{-1}(c)} \mathrm{pop}(v)$.
The initial coarse plan $\pi_C^{(0)} : C \to [k]$ is derived from the
initial fine plan $\pi^{(0)} : V \to [k]$ by \emph{plurality projection}:
\[
  \pi_C^{(0)}(c) = \arg\max_{d \in [k]}\;
    \mathrm{pop}\!\left(\{v \in f^{-1}(c) : \pi^{(0)}(v) = d\}\right).
\]
That is, each county is assigned to the district that contains the plurality
of its population.

\subsection{Rebalance Subroutine}

After a coarse move proposes a new $\pi_C'$, the fine plan is updated by
assigning each tract to the district of its county:
$\pi'(v) = \pi_C'(\mathrm{county}(v))$.
This projection may violate fine population balance because counties are not
perfectly equal in population.
The \emph{rebalance subroutine} restores balance by performing boundary swaps:

\begin{enumerate}
  \item Identify all districts $d$ with $|\mathrm{pop}(d) - P_{\mathrm{ideal}}|
        > \epsilon_f \cdot P_{\mathrm{ideal}}$.
  \item For each over-populated district, find boundary tracts adjacent to
        under-populated districts and move the tract with the best
        population-balance improvement.
  \item Repeat up to 200 iterations.
  \item If any district remains unbalanced after 200 iterations, \emph{reject}
        the coarse move entirely (revert to the previous fine plan $\pi$).
\end{enumerate}

The 200-iteration limit is a hard contract: it prevents the rebalance from
consuming unbounded time while still accommodating the typical imbalances
introduced by county-level projection.
In practice, for Texas, rebalance completes within 5--20 iterations on
accepted coarse moves; failure occurs when a coarse move produces a coarse
district that contains an entire county whose population far exceeds
$P_{\mathrm{ideal}}$.

\subsection{Complete Algorithm}

\begin{algorithm}
\caption{\ms{} Option B}
\label{alg:multiscale}
\begin{algorithmic}[1]
\Require Tract graph $G$, populations $\{\mathrm{pop}(v)\}$, GEOIDs,
         initial plan $\pi^{(0)}$, districts $k$,
         fine tolerance $\epsilon_f$, coarse factor 3,
         coarsening probability $\alpha$, total steps $T$,
         percentile $p \in [0,1]$, base seed $b$
\Ensure  Selected plan $\pi^*$

\State $\epsilon_c \gets 3 \epsilon_f$
\State $f(v) \gets \mathrm{geoid}(v)[0{:}5]$ for all $v \in V$
        \Comment{tract $\to$ county map}
\State $G_C, \{\mathrm{pop}(c)\} \gets \textsc{DeriveCountyGraph}(G, f)$
        \Comment{$O(|E|)$, Prop.~\ref{prop:county-adj-exact}}
\State $\pi_C \gets \textsc{PluralityProject}(\pi^{(0)}, f)$
\State $\pi \gets \pi^{(0)}$
\State $\mathrm{accepted} \gets [(\ec(\pi),\, \pi)]$

\For{$s \gets 1$ \textbf{to} $T$}
  \State $u \gets \mathrm{SHA256}(\texttt{"MSC\_STEP\_"} \,\|\, s^{\mathrm{u64le}}
         \,\|\, 0^{\mathrm{u32le}} \,\|\, b^{\mathrm{u64le}})[0{:}8]$
         as $\mathrm{SmallRng}$
  \If{$u.\mathrm{gen\_range}(0.0, 1.0) < \alpha$}
    \Comment{Coarse move}
    \State $\pi_C' \gets \recom{}.\mathrm{step}(G_C, \{\mathrm{pop}(c)\}, \pi_C, \epsilon_c, u)$
    \State $\pi' \gets \textsc{Project}(\pi_C', f)$
            \Comment{assign each tract to its county's district}
    \State $\mathrm{ok} \gets \textsc{Rebalance}(\pi', \epsilon_f, \text{max}=200)$
    \If{$\mathrm{ok}$}
      \State $\pi \gets \pi'$;\quad $\pi_C \gets \pi_C'$
      \State $\mathrm{accepted}.\mathrm{push}((\ec(\pi), \pi))$
    \EndIf
    \Comment{else reject: keep $\pi$, $\pi_C$ unchanged}
  \Else
    \Comment{Fine move}
    \State $\pi \gets \recom{}.\mathrm{step}(G, \{\mathrm{pop}(v)\}, \pi, \epsilon_f, u)$
    \State $\pi_C \gets \textsc{PluralityProject}(\pi, f)$
            \Comment{keep coarse plan consistent}
    \State $\mathrm{accepted}.\mathrm{push}((\ec(\pi), \pi))$
  \EndIf
\EndFor

\State Sort $\mathrm{accepted}$ by $\ec$ ascending
\State \Return $\mathrm{accepted}[\lfloor p \cdot |\mathrm{accepted}|\rfloor]$
\end{algorithmic}
\end{algorithm}

\paragraph{Deterministic seed schedule.}
Each step $s$ derives its RNG from a SHA-256 hash of the step index and
the base seed:
\[
  \mathrm{step\_seed}(s, b)
  = \mathrm{SHA256}\!\left(\texttt{"MSC\_STEP\_"}
    \,\|\, s^{\mathrm{u64le}}
    \,\|\, 0^{\mathrm{u32le}}
    \,\|\, b^{\mathrm{u64le}}\right)[0{:}8]
    \;\text{as}\; u64.
\]
This ensures that the sequence of coarse/fine decisions---and the outcomes
of each step---is fully determined by $b$ and $T$, enabling exact
reproduction of any run from the manifest entry alone.

\paragraph{CLI invocation.}
\begin{verbatim}
redist state --state TX --year 2020 \
    --search multiscale \
    --multiscale-steps 5000 \
    --multiscale-alpha 0.3
\end{verbatim}

\subsection{Correctness Properties}

\begin{proposition}[Fine plan validity]
\label{prop:validity}
Every fine plan $\pi$ in the \texttt{accepted} list is connected and
population-balanced within $\epsilon_f$.
\end{proposition}

\begin{proof}
The initial plan $\pi^{(0)}$ is valid by assumption.
Fine \recom{} steps preserve validity (the standard \recom{} invariant).
Coarse moves are accepted only if \textsc{Rebalance} returns \texttt{ok},
which certifies $|\mathrm{pop}(d) - P_{\mathrm{ideal}}| \leq \epsilon_f \cdot
P_{\mathrm{ideal}}$ for all $d$ after rebalancing.
Connectivity of fine districts after a coarse move: each tract is assigned
to the district of its county; within a county, the district assignment is
uniform, so connectivity of the county in $G_C$ is necessary but not
sufficient for connectivity of the fine district.
The rebalance swaps maintain a connected fine district by restricting to
boundary-tract swaps that do not disconnect the donor district (checked via
BFS, as in the \flip{} chain).
\end{proof}

\begin{remark}[Stationarity]
\label{rem:stationarity}
Like standard \recom{}, \ms{} does not have a fully characterised stationary
distribution.
The coarse moves introduce an additional source of non-reversibility: the
projection and rebalance operations are not symmetric in forward and reverse
directions.
For practical redistricting ensemble analysis, this is acceptable---the goal
is faster mixing to explore the plan space, not exact sampling from a specific
distribution.
Applications requiring exact distributional guarantees should use
\fr{}~\citep{dellaLibera2026forestRecom} at the fine level, though the
multi-scale extension of \fr{} is left for future work.
\end{remark}
